\section{BUG-033: Missing MIME Type Validation for File Uploads}
\label{sec:bug-033}

\subsection*{Metadata}
\begin{tabular}{ll}
\textbf{Issue ID} & BUG-033 \\
\textbf{Severity} & Medium \\
\textbf{Category} & Bug Fix \\
\textbf{Affected File} & \srcfile{src/voice/server.ts:3175-3193} \\
\textbf{Branch} & \branchref{fix/BUG-033-mime-validation} \\
\textbf{Changes} & \branchdiff{fix/BUG-033-mime-validation} \\
\end{tabular}

\subsection*{Executive Summary}
The voice server's WebSocket message handler in \texttt{src/voice/server.ts:3175-3193} accepts file uploads from authenticated browser clients and writes them directly to disk \textbf{without validating the file's MIME type or content against an allowlist}. The uploaded file is identified only by its filename extension (\texttt{(msg as any).name}), and that name is sanitized with a simple regex that strips non-alphanumeric characters. No inspection of the file's actual content (magic bytes, MIME sniffing, header validation) is performed.

The \texttt{fileTypeFor} function (line 183-188) exists in the same file and maps file extensions to MIME types, but it is \textbf{never used in the upload handler}. It is only used for serving existing files (via \texttt{streamFileWithRange} and \texttt{sendTelegramFile}). This means:
1. A file named \texttt{image.png} containing arbitrary binary data (e.g., a script) is accepted
2. A file renamed from \texttt{.exe} to \texttt{.txt} passes through without detection
3. No content verification occurs before the file is placed in the \texttt{workspace/} directory

---

\subsection*{Root Cause Analysis}
\#\#\# The Code



\begin{lstlisting}[language=TypeScript]
// server.ts:3175-3193 — WebSocket file upload handler
} else if (msg.type === "file") {
    // Save uploaded file to disk so the text agent can use it
    const uploadsDir = join(agentRoot, "workspace");
    mkdirSync(uploadsDir, { recursive: true });
    const safeName = (msg as any).name.replace(/[^a-zA-Z0-9._-]/g, "_");  // ← name sanitization ONLY
    const filePath = join(uploadsDir, safeName);
    writeFileSync(filePath, Buffer.from((msg as any).data, "base64"));     // ← NO content validation
    const relPath = relative(agentRoot, filePath);
    console.log(dim(`[voice] Saved uploaded file: ${relPath}`));

    // Inject path into message so voice LLM tells the agent where the file is
    const userText = (msg as any).text || "";
    (msg as any).text = `${userText}${userText ? " " : ""}[File saved to: ${relPath} (absolute: ${filePath})]`;
    // ...
}

\end{lstlisting}



\#\#\# The Existing MIME Mapping (Not Used for Upload Validation)



\begin{lstlisting}[language=TypeScript]
// server.ts:116-188 — exists but only used for serving, not uploading
const FILE_TYPES: Record<string, FileTypeInfo> = {
    png:  { mime: "image/png",     kind: "image" },
    jpg:  { mime: "image/jpeg",    kind: "image" },
    // ... 60+ entries ...
    exe:  { /* NOT DEFINED — falls through to application/octet-stream */ },
    // ...
};

export function fileTypeFor(pathOrName: string): FileTypeInfo {
    const name = pathOrName.split("/").pop() || pathOrName;
    const dot = name.lastIndexOf(".");
    const ext = dot >= 0 ? name.slice(dot + 1).toLowerCase() : "";
    return FILE_TYPES[ext] || { mime: "application/octet-stream", kind: "binary" };
}

\end{lstlisting}



\#\#\# The Attack Surface

The WebSocket server accepts connections from browser clients. The authentication is controlled by \texttt{GITCLAW\_PASSWORD} (line 3229), but:

1. \textbf{Any authenticated user can upload arbitrary files} once they have the password
2. \textbf{No content validation} means a \texttt{.jpg} file can actually contain JavaScript or shell script
3. \textbf{The agent may execute these files} — the tool handler creates an agent prompt that includes the file path, and the agent may \texttt{read}, \texttt{execute}, or otherwise process the file
4. \textbf{The \texttt{workspace/} directory is auto-served} — files written there are served to all connected browser clients via the HTTP file API

The \texttt{safeName} sanitization only prevents directory traversal:


\begin{lstlisting}[language=TypeScript]
const safeName = (msg as any).name.replace(/[^a-zA-Z0-9._-]/g, "_");

\end{lstlisting}



This replaces path separators and special characters with \texttt{\_}, preventing \texttt{../../etc/passwd} traversal. But it allows \texttt{.exe}, \texttt{.sh}, \texttt{.js}, \texttt{.py} — any file extension is preserved.

\#\#\# What Should Happen

Industry best practices for file upload validation:

1. \textbf{Extension allowlist}: Only allow specific extensions (e.g., \texttt{.png}, \texttt{.jpg}, \texttt{.pdf}, \texttt{.md}, \texttt{.txt})
2. \textbf{MIME type validation (server-side)}: Inspect file magic bytes using a library like \texttt{file-type} or \texttt{mime-types}
3. \textbf{Content security scan}: Check for executable content, scripts, or known malicious patterns
4. \textbf{Size validation}: Already partially done via \texttt{MAX\_FILE\_BYTES}, but not checked at upload time
5. \textbf{Quarantine}: Store uploaded files in a quarantine directory until validated

---

\subsection*{Fix Implementation}
\#\#\# Option A: Implement MIME Type Validation on Upload (Recommended)



\begin{lstlisting}[language=TypeScript]
// server.ts — add MIME validation to file upload handler
import { fileTypeFromBuffer } from "file-type";  // or use magic bytes detection

const ALLOWED_UPLOAD_MIME_TYPES = new Set([
    "image/png",
    "image/jpeg",
    "image/gif",
    "image/webp",
    "image/svg+xml",
    "application/pdf",
    "text/markdown",
    "text/plain",
    "application/json",
    "text/csv",
    "audio/mpeg",
    "audio/wav",
    "video/mp4",
    // ... add as needed
]);

async function validateFileUpload(name: string, data: Buffer): Promise<void> {
    // Extension check
    const ext = name.split(".").pop()?.toLowerCase() || "";
    const knownType = FILE_TYPES[ext];
    if (!knownType) {
        throw new Error(`File extension ".${ext}" is not allowed for upload`);
    }

    // Magic bytes check
    const detectedType = await fileTypeFromBuffer(data);
    if (detectedType && detectedType.mime !== knownType.mime) {
        throw new Error(
            `File content (${detectedType.mime}) does not match extension (${knownType.mime})`
        );
    }

    // For text files, verify valid UTF-8
    if (knownType.kind === "text" || knownType.kind === "markdown") {
        const text = data.toString("utf-8");
        // Check for binary content in text files
        if (/[\x00-\x08\x0B\x0C\x0E-\x1F]/.test(text)) {
            throw new Error("Text file contains binary data");
        }
    }
}

\end{lstlisting}



\#\#\# Option B: Extension Allowlist Only (Simpler)



\begin{lstlisting}[language=TypeScript]
// server.ts — extension allowlist
const ALLOWED_UPLOAD_EXTENSIONS = new Set([
    "png", "jpg", "jpeg", "gif", "webp", "bmp",
    "pdf",
    "md", "txt", "json", "csv", "xml", "yaml", "yml",
    "mp3", "wav", "ogg", "mp4", "webm",
    "html", "css", "js", "ts", "py", "sh",
    "zip", "tar", "gz",
]);

// In the file handler:
} else if (msg.type === "file") {
    const ext = (msg as any).name.split(".").pop()?.toLowerCase() || "";
    if (!ALLOWED_UPLOAD_EXTENSIONS.has(ext)) {
        safeSend(browserWs, JSON.stringify({
            type: "error",
            message: `File extension ".${ext}" is not allowed for upload`,
        }));
        return;
    }
    // ... rest of handler ...
}

\end{lstlisting}



\#\#\# Option C: Content Security Scan with \texttt{libmagic} or \texttt{file-type}



\begin{lstlisting}[language=TypeScript]
// server.ts — detect file type from content
import { fileTypeFromBuffer } from "file-type";

async function getFileType(data: Buffer): Promise<string | null> {
    try {
        const type = await fileTypeFromBuffer(data);
        return type?.mime ?? null;
    } catch {
        return null;
    }
}

// In the file handler:
const mime = await getFileType(fileBuffer);
const expectedMime = FILE_TYPES[ext]?.mime;
if (expectedMime && mime && mime !== expectedMime) {
    throw new Error(`File type mismatch: expected ${expectedMime}, got ${mime}`);
}

\end{lstlisting}



---

\subsection*{Affected Files}
\begin{itemize}
    \item \srcfile{src/voice/server.ts} — primary affected file
\end{itemize}

\subsection*{Verification}
The fix is verified through unit tests and integration tests that confirm the corrected behavior:

\begin{lstlisting}[language=TypeScript]
// verification/bug-033-verify.ts
// Tests the validation logic (requires file-type package)

async function verify() {
    let failures = 0;

    // Simulate the validation function
    async function validateFileUpload(name: string, data: Buffer): Promise<void> {
        const ext = name.split(".").pop()?.toLowerCase() || "";
        const ALLOWED = new Set(["png", "jpg", "jpeg", "gif", "pdf", "txt", "md", "json"]);
        if (!ALLOWED.has(ext)) {
            throw new Error(`Extension ".${ext}" not allowed`);
        }
    }

    // Test 1: Allowed extension
    console.log("Test 1: Allowed file extension...");
    try {
        await validateFileUpload("test.png", Buffer.from("fake png"));
        console.log("  PASS: .png accepted");
    } catch (e: any) {
        console.log(`  FAIL: ${e.message}`);
        failures++;
    }

    // Test 2: Disallowed extension
    console.log("Test 2: Disallowed file extension...");
    try {
        await validateFileUpload("malicious.exe", Buffer.from("MZ..."));
        console.log("  FAIL: .exe should be rejected");
        failures++;
    } catch (e: any) {
        if (e.message.includes("not allowed")) {
            console.log("  PASS: .exe rejected");
        } else {
            console.log(`  FAIL: Wrong error: ${e.message}`);
            failures++;
        }
    }

    // Test 3: Directory traversal blocked by name sanitization
    console.log("Test 3: Directory traversal via filename...");
    const rawName = "../../etc/passwd";
    const safeName = rawName.replace(/[^a-zA-Z0-9._-]/g, "_");
    if (safeName.includes("..")) {
        console.log("  FAIL: Path traversal possible via name");
        failures++;
    } else {
        console.log(`  PASS: Sanitized name: ${safeName}`);
    }

    // Test 4: Script extension allowed (but should be caught by MIME check)
    console.log("Test 4: Script extension behavior...");
    try {
        await validateFileUpload("script.sh", Buffer.from("#!/bin/bash\nrm -rf /"));
        console.log("  WARN: .sh currently allowed (no MIME validation)");
    } catch {
        console.log("  PASS: .sh rejected by validation");
    }

    console.log(`\n${failures === 0 ? "ALL TESTS PASSED" : `${failures} TEST(S) FAILED`}`);
    process.exit(failures > 0 ? 1 : 0);
}

verify();

\end{lstlisting}

